\section{Training Strategies for Speech LLMs}
\label{sec:training_strategies}

Building Speech LLMs requires extending the text-LLM training pipeline to
account for the continuous, multi-channel nature of speech. We organize
the literature around four training stages: (i) modality alignment and
pretraining, (ii) fine-tuning and adaptation, (iii) preference alignment,
and (iv) cross-cutting design issues that connect training to the
remainder of the lifecycle. Throughout, we emphasize choices that have
downstream implications for evaluation
(Section~\ref{sec:evaluation_benchmarking}), deployment
(Section~\ref{sec:deployment}), and safety
(Section~\ref{sec:safety_privacy_governance}).

\subsection{Modality Alignment and Pretraining}
\label{subsec:training_pretraining}

Speech and text differ in two ways that complicate joint modeling: speech
sequences are typically much longer than their text counterparts, and
speech embeddings live in a separate representation space from text
embeddings \citep{fan2025alignformer,bai2024seed}. Modality alignment
modules and acoustic pretraining strategies address these mismatches.

\paragraph{Alignment modules.}
Modular Speech LLMs combine a pretrained speech encoder (e.g., Whisper
\citep{radford2023robust}, HuBERT \citep{hsu2021hubert}, wav2vec~2.0
\citep{baevski2020wav2vec}) with a partly frozen text LLM through a
trainable bridge. \emph{AlignFormer} \citep{fan2025alignformer} pairs a
Connectionist Temporal Classification (CTC) layer with a Q-Former-style
attention module to map variable-length acoustic features to discrete
text-aligned tokens. \emph{Seed-ASR Converter} \citep{bai2024seed} learns
a context-aware transformation that improves multilingual generalization.
\emph{SALMONN}'s Q-Former \citep{tang2023salmonn} and the OSUM adapter
\citep{geng2025osum} use a small set of trainable queries to distill
semantically rich acoustic representations into the LLM's input space.
As an alternative to adapter-based bridging, X-OPD
\citep{cao2026xopd} distills text-LLM capabilities into a speech-LLM via
cross-modal on-policy rollouts, bypassing the need for paired
speech-text supervision altogether. Empirical alignment quality can be
probed via the Automatic Latent Alignment Score (ALAS)
\citep{mousavi2025alas}, which uses Whisper timestamps as a reference and
reports that deeper transformer layers align better for semantic tasks
while earlier layers align better for paralinguistic tasks.

\paragraph{Self-supervised pretraining.}
Speech LLMs adapt the masked-prediction and clustering objectives from
text SSL to raw audio. Representative methods include UniSpeech-SAT
\citep{chen2022unispeech}, which adds speaker-aware objectives to the SSL
loss; DinoSR \citep{liu2023dinosr}, which combines self-distillation with
masked acoustic modeling; and Seed-ASR \citep{bai2024seed}, which scales a
BERT-style masked-prediction pretraining to millions of hours of speech
to produce the LUISE encoder, subsequently aligned to a language model.

\paragraph{Joint and weakly supervised pretraining.}
Two further pretraining paradigms reduce reliance on parallel speech-text
data. \emph{Cold-start joint modeling} initializes both speech and
language components together, as in Generative Spoken Language Modeling
(GSLM) \citep{lakhotia2021generative} and SPIRIT-LM \citep{nguyen2025spirit},
yielding tighter cross-modal interactions than methods that freeze the
language model. \emph{Weakly supervised pretraining at scale}, exemplified
by Whisper \citep{radford2023robust}'s 680{,}000-hour multilingual corpus,
trades exact transcription quality for breadth and produces strong
zero-shot transcription, translation, and detection capabilities. Many
later systems (OSUM \citep{geng2025osum}, SALMONN \citep{tang2023salmonn})
build directly on Whisper-style encoders, treating large-scale weak
supervision as the de facto foundation.

\subsection{Fine-Tuning and Adaptation}
\label{subsec:training_finetuning}

Pretraining yields foundations; fine-tuning produces task-aligned
behavior. Speech LLM fine-tuning takes three main forms.

\paragraph{Instruction tuning.}
Instruction tuning aligns Speech LLMs to user intent in spoken dialogue,
question answering, and editing tasks. SpeechGPT
\citep{zhang2023speechgptempoweringlargelanguage} integrates discrete
speech tokens into a decoder-only LLM backbone with strong zero-shot
generalization. COSMIC \citep{pan2024cosmicdataefficientinstructiontuning}
trains a Q-Former bridge on 450 hours of speech plus GPT-3.5-generated
QA pairs, achieving competitive ASR and speech-to-text translation with
limited paired data. InstructSpeech \citep{huang2024instructspeech} uses
hierarchical adapters and multi-step reasoning for fine-grained semantic
and acoustic editing. LLaMA-Omni
\citep{fang2025llamaomniseamlessspeechinteraction} pairs a frozen Whisper
encoder with LLaMA-3.1-8B-Instruct and a streaming decoder, reporting
226\,ms response latency without intermediate transcription. More
recently, Audio-FLAN \citep{xue2025audioflan} demonstrates that scaling
instruction data to over 100 million instances across 80 audio tasks
yields substantial gains, signaling a shift toward
massive-scale instruction tuning.

\paragraph{Parameter-efficient tuning.}
Parameter-efficient fine-tuning (PEFT) updates a small fraction of model
parameters while freezing the rest, reducing compute cost and mitigating
catastrophic forgetting. HDMoLE \citep{mu2025hdmolemixtureloraexperts}
combines hierarchical LoRA adapters with mixture-of-experts routing for
joint accent and domain adaptation, updating $\sim$9.6\% of parameters
while matching full fine-tuning. ELP-Adapters
\citep{inoue2024elpadaptersparameterefficientadapter} insert lightweight
modules at multiple abstraction levels of pretrained encoders such as
WavLM, enabling targeted adaptation with minimal parameter footprint.
Zipper-LoRA \citep{mei2026zipperlora} further advances speech-specific
PEFT by introducing language-conditioned routing with rank-level
decoupling for multilingual speech recognition.

\paragraph{Domain-specific tuning.}
Specialized domains (medical, financial, technical) often have
constrained paired speech-text resources. Low-resource adaptation via
unpaired domain-specific text \citep{fang2025lowresourcedomainadaptationspeech}
fine-tunes the language decoder with LoRA while preserving alignment via
real-time evaluation, reducing domain WER without harming general
performance. Decoder-only context perturbation \citep{suh2024improving}
freezes the encoder and randomizes domain descriptions during decoder
fine-tuning to mitigate overfitting and forgetting, with successful
applications to earnings calls and medical dialogue.

\subsection{Post-Training Preference Alignment}
\label{subsec:training_alignment}

Pretraining and fine-tuning equip Speech LLMs with capability; preference
alignment shapes their behavior to match user expectations for
helpfulness, safety, and naturalness. This is especially consequential in
speech, where prosody, pacing, and emotion strongly influence perceived
quality \citep{communication2023seamlessmultilingualexpressivestreaming}.

\paragraph{Reinforcement Learning from Human Feedback (RLHF).}
RLHF \citep{ouyang2022training,christiano2017deep} trains a reward model
on human-labeled preference pairs and then optimizes the base model
against it. The reward model uses a pairwise ranking loss that converts
human comparisons into a scalar reward:
\begin{equation}
\mathcal{L}^{\text{reward}}(\phi) = -\log \sigma\left( r_\phi(x, y_1) - r_\phi(x, y_2) \right),
\end{equation}
where $x$ is the prompt, $y_1$ is preferred over $y_2$, and $r_\phi$
scores responses. The base model is then optimized via policy-gradient
methods (typically PPO \citep{schulman2017proximalpolicyoptimizationalgorithms}
in text settings) with a KL constraint to a reference model. RLHF is
expensive in human annotation and unstable to tune, particularly in
speech, where preference data must capture acoustic and paralinguistic
quality in addition to content. As a result, Speech LLMs typically
substitute RLHF with lighter alternatives \citep{rafailov2023direct}.

\paragraph{Direct Preference Optimization (DPO).}
DPO \citep{rafailov2023direct} eliminates the explicit reward model by
expressing the preference loss directly in terms of the policy:
\begin{equation}
\mathcal{L}_{\mathrm{DPO}}(\theta)
= - \mathbb{E}_{x,y_w,y_l}
\left[
\log \sigma \left(
\beta \left(
\log \frac{\pi_\theta(y_w \mid x)}{\pi_{\mathrm{ref}}(y_w \mid x)}
-
\log \frac{\pi_\theta(y_l \mid x)}{\pi_{\mathrm{ref}}(y_l \mid x)}
\right)
\right)
\right].
\end{equation}

DPO matches or exceeds RLHF on text alignment tasks while using a single
training stage. Speech LLMs have rapidly adopted it. SpeechAlign
\citep{zhang2024speechalignaligningspeechgeneration} iteratively refines a
zero-shot TTS model on LibriSpeech and VCTK, surpassing RLHF-PPO on Word
Error Rate, Speaker Similarity, and human ratings. Align-SLM
\citep{lin2025alignslmtextlessspokenlanguage} extends DPO to a textless
spoken language model trained on LibriSpeech and Multilingual LibriSpeech,
outperforming Moshi and SPIRIT-LM on the spoken StoryCloze benchmark.
OpenOmni \citep{luo2025openomniadvancingopensourceomnimodal} introduces a
Connectionist Temporal Classification variant of DPO (CTC-DPO) for
emotionally coherent speech generation, reporting strong results on the
EO2S-9k bilingual emotional speech benchmark. Beyond pairwise
comparisons, Emo-LiPO \citep{lin2026emolipo} proposes listwise preference
optimization for speech emotion intensity, suggesting that richer
ranking signals can further improve alignment quality.

\paragraph{Reinforcement Learning from AI Feedback (RLAIF).}
RLAIF \citep{lee2024rlaifvsrlhfscaling,bai2022constitutionalaiharmlessnessai}
replaces human preference annotation with AI-generated labels, trading a
small accuracy cost for substantially improved scalability. Lee et al.\
show that direct RLAIF (d-RLAIF) prompts an AI labeler at training time
and matches or exceeds RLHF on summarization and dialogue. RLAIF-V
\citep{yu2024rlaifvopensourceaifeedback} extends this to multimodal models
with deconfounded generation and atomic claim decomposition; the 12B
variant reduces its overall hallucination rate by 32.4\% through
self-alignment, surpassing GPT-4V in trustworthiness. In speech,
AlignSLM \citep{lin2025alignslmtextlessspokenlanguage} uses semantic
metrics on AI-generated speech continuations to construct preference data
for DPO, achieving state-of-the-art results on spoken StoryCloze.

\paragraph{Comparison.}
Table~\ref{tab:alignment_comparison} compares RLHF, DPO, and RLAIF
alongside their speech-specific evaluation considerations. The choice
depends not only on annotation cost, but also on what evaluation regime
(Section~\ref{sec:evaluation_benchmarking}) is available: methods that
require dense per-turn or per-utterance ratings are difficult when
benchmarks can only score complete responses.

\begin{table*}[t]
\caption{Comparison of preference alignment methods for Speech LLMs. The
last column links each method to evaluation considerations developed in
Section~\ref{sec:evaluation_benchmarking}.}
\label{tab:alignment_comparison}
\centering
\scriptsize
\begin{tabularx}{\linewidth}{lXXXX}
\toprule
\textbf{Aspect} & \textbf{RLHF} & \textbf{DPO} & \textbf{RLAIF} & \textbf{Speech-evaluation hook} \\
\midrule
\textbf{Annotation source} & Human preference pairs & Human preference pairs; some implementations mix synthetic data & AI-generated preference pairs & Speech requires acoustic + paralinguistic judgements; LLM judges are still moderate-correlation \citep{zheng2023judging} \\
\midrule
\textbf{Optimization} & Policy gradient (e.g., PPO) & Direct preference loss & Reward learning or DPO-style & Construct/evaluator validity (Section~\ref{sec:evaluation_benchmarking}) determines which optimization signal is reliable \\
\midrule
\textbf{Reward model} & Required & Not required & Required (or implicit) & A reward trained on transcripts misses acoustic quality; speech reward models remain underdeveloped \\
\midrule
\textbf{Engineering overhead} & High (multi-stage pipeline) & Moderate (single-stage) & Moderate (auto-labeling pipeline) & Operational validity hooks: alignment cost shapes deployment cycle (Section~\ref{sec:deployment}) \\
\midrule
\textbf{Speech adoption} & Limited (manual effort, slow iteration) & Promising: SpeechAlign, Align-SLM, OpenOmni & Emerging: RLAIF-V, AlignSLM & Adoption tracks the availability of speech-aware judges and benchmarks \\
\midrule
\textbf{Scalability} & Limited by annotation cost & Improved efficiency & High (AI-supervised) & Scaling depends on whether evaluation can detect over-optimization (Section~\ref{sec:evaluation_benchmarking}) \\
\bottomrule
\end{tabularx}
\end{table*}

\subsection{Cross-Cutting Training Issues}
\label{subsec:training_crosscutting}

Three training-design decisions have outsized influence on the rest of
the lifecycle and motivate later sections.

\emph{Data sufficiency.} Self-supervised and weakly supervised
pretraining on tens to hundreds of thousands of hours has become the
default \citep{radford2023robust,bai2024seed}, but parallel speech-text
data remains scarce for low-resource languages, scripts, and domains.
This data imbalance reappears in evaluation: most multilingual
benchmarks under-sample non-Latin scripts (Section~\ref{sec:evaluation_benchmarking}).

\emph{Alignment quality.} Whether alignment modules preserve
paralinguistic content varies substantially across architectures
\citep{mousavi2025alas,fan2025alignformer}. Benchmarks such as SALMon and
FMSU-Bench show that semantically capable Speech LLMs frequently fail at
acoustic-semantic binding (Section~\ref{sec:evaluation_benchmarking}),
which suggests that alignment quality, not language modeling, is often
the bottleneck.

\emph{Deployment-aware training.} Few training pipelines optimize
explicitly for streaming, latency, or interruption handling. Models
designed for offline benchmarks may compute correctly but respond too
late for natural interaction, motivating the architecture and protocol
choices analyzed in Section~\ref{sec:deployment}. Safety-aware training
is similarly under-developed: text safety alignment does not
automatically transfer to audio inputs, as
Section~\ref{sec:safety_privacy_governance} demonstrates.

These cross-cutting issues connect the training section to the rest of the survey. The following sections examine whether current evaluation, deployment, and safety practices can measure the limitations introduced by training choices.
